\documentclass[11pt,a4paper]{coverletter}

\senderblock
  {Radheem Bin Razi}
  {Distributed Systems \& Cloud-Native Engineer}
  {Ilmenau, Germany \textbar\ \href{mailto:sheikh.radheem@gmail.com}{sheikh.radheem@gmail.com} \textbar\ \href{https://radheem.github.io/my-cv}{portfolio}}

\begin{document}

%% ===== ENGLISH =====
\recipient{Bundeswehr (German Armed Forces)}{Hiring Team}{}

\opening{Dear Hiring Team,}

\vspace{8pt}\noindent\textbf{1. Why Bundeswehr?}\par\vspace{3pt}

Designing a resilient private cloud storage platform requires rigorous automation and deep infrastructure fundamentals. I am drawn to the Bundeswehr’s commitment to Zero-Touch operations and long-term cloud strategy. This aligns with my professional goals to apply my Kubernetes, GitOps, and Linux automation background to secure, scalable systems. While Rook and Ceph are not yet in my daily toolkit, my experience with Cilium networking, bare-metal cluster orchestration, and infrastructure-as-code ensures a rapid ramp-up. I would like to contribute my systems engineering experience to help your team build a reliable storage foundation.

\vspace{8pt}\noindent\textbf{2. Why me?}\par\vspace{3pt}

\bullets

  \item Took end-to-end ownership of a Go microservices platform, shipping production-ready APIs, NATS JetStream messaging, and Kubernetes deployments using kustomize and Skaffold.

  \item Collaborated across product, security, and infrastructure teams to integrate third-party vendors, establish CI/CD pipelines, and enforce consistent schema validation via CUE.

  \item Engineered a reboot-resilient K3s cluster with Cilium eBPF networking, cert-manager, external-dns, and a sequential-boot orchestrator that self-heals CNI failures without kube-proxy.

\bulletsend

I am available full-time immediately, can relocate within Germany within two to three weeks, hold Blue Card readiness so your team only needs to issue the contract, and maintain an A2 German level that I am actively advancing.

\closing{Sincerely,}{Radheem Bin Razi}

\clearpage
\selectlanguage{ngerman}

%% ===== DEUTSCH =====
\recipient{Bundeswehr (German Armed Forces)}{Personalabteilung}{}

\opening{Sehr geehrtes Hiring-Team,}

\vspace{8pt}\noindent\textbf{1. Warum Bundeswehr?}\par\vspace{3pt}

Der Aufbau einer resilienten privaten Cloud-Speicherplattform erfordert strikte Automatisierung und fundierte Infrastrukturkenntnisse. Das Engagement der Bundeswehr für Zero-Touch-Betrieb und eine langfristige Cloud-Strategie spricht mich besonders an. Dies entspricht meinen beruflichen Zielen, mein Hintergrundwissen in Kubernetes, GitOps und Linux-Automatisierung für sichere und skalierbare Systeme einzusetzen. Auch wenn Rook und Ceph noch nicht zu meinem täglichen Werkzeugkasten gehören, gewährleisten meine Erfahrungen mit Cilium-Netzwerken, der Orchestrierung von Bare-Metal-Clustern und Infrastructure-as-Code eine rasante Einarbeitungszeit. Gerne möchte ich meine Erfahrung im Systems Engineering einbringen, um Ihr Team beim Aufbau einer zuverlässigen Speicherbasis zu unterstützen.

\vspace{8pt}\noindent\textbf{2. Warum ich?}\par\vspace{3pt}

\bullets

  \item Übernahm die End-to-End-Verantwortung für eine Go-Mikroservices-Plattform, lieferte produktionsreife APIs, NATS JetStream-Messaging sowie Kubernetes-Deployments unter Verwendung von kustomize und Skaffold.

  \item Zusammenarbeit mit Produkt-, Sicherheits- und Infrastrukturteams zur Integration von Drittanbieterlösungen, Aufbau von CI/CD-Pipelines sowie Durchsetzung einer konsistenten Schema-Validierung mittels CUE.

  \item Aufbau eines neustartresistenten K3s-Clusters mit Cilium eBPF-Netzwerken, cert-manager, external-dns sowie einem sequenziellen Boot-Orchestrator, der CNI-Ausfälle ohne kube-proxy selbstständig repariert.

\bulletsend

Ich bin ab sofort vollzeit verfügbar, kann innerhalb von zwei bis drei Wochen innerhalb Deutschlands umziehen, bin für die Blue Card vorbereitet, sodass Ihr Team lediglich den Arbeitsvertrag ausstellen muss, und verfüge über Deutschkenntnisse auf A2-Niveau, die ich aktiv ausbaue.

\closing{Mit freundlichen Grüßen,}{Radheem Bin Razi}

\end{document}
